\section{Introduction}\label{sec_intro}

Let $G=\mathrm{GL}_n(\C)$, let $B\subset G$ be the Borel subgroup of upper triangular matrices, let $T\subset B$ be the diagonal maximal torus, and write $\CB=G/B$ for the complete flag variety. Let $\mathfrak g=\mathfrak{gl}_n(\C)$ be the Lie algebra of $G$, and let $\Delta^+$ denote the set of positive roots of $(G,T)$ corresponding to $B$, with root spaces $\mathfrak g_\alpha\subset\mathfrak g$. Via the trace pairing on $\mathfrak g$, the cotangent bundle $T^*\CB$ is identified with $G\times^B\bigoplus_{\alpha\in\Delta^+}\mathfrak g_\alpha$. Associated with each root ideal $\Psi\subset\Delta^+$, Kato \cite[Theorems~3.10 and 4.4]{kato2023geometric} constructed a smooth projective compactification $\SX_\Psi$ of the $G$-equivariant vector subbundle
\[
 T^*_\Psi\CB=G\times^B\bigoplus_{\alpha\in\Psi}\mathfrak g_\alpha\subset T^*\CB.
\]
We refer to $\SX_\Psi$ as a \emph{Kato Catalan variety}.\footnote{In \cite{kato2023geometric}, Kato did not assign a name to $\SX_\Psi$; in private communication, he mentioned that he had originally intended to call it a \emph{Catalan variety}, but refrained from doing so as that terminology was already in use elsewhere.} These varieties were introduced to provide a geometric realization of Catalan functions and, via the endpoint morphism to the affine Grassmannian $\Gr$ and the geometric Satake equivalence, of chromatic quasisymmetric polynomials \cite{kato2023geometric,kato2024geometric,kamnitzer2024fixed}. From the viewpoint of $\Gr$, Kamnitzer described $\SX_\Psi$ as a reverse Hessenberg convolution variety and exhibited it as a tower of projective bundles over a point \cite[Proposition~2.1 and Remark~2.3]{kamnitzer2024fixed}.

For the maximal root ideal $\Psi=\Delta^+$, let $\CN\subset\mathfrak{gl}_n(\C)$ denote the nilpotent cone. The composite of the open embedding $T^*\CB\hookrightarrow\SX_{\Delta^+}$ with the endpoint morphism $\mathsf m_{\Delta^+}:\SX_{\Delta^+}\to\Gr$ is the composite of the Springer resolution $T^*\CB\to\CN$ with the classical lattice embedding $\CN\hookrightarrow\Gr$. This lattice construction originates with Lusztig \cite{lusztig1981green} and was developed further by Ng\^o \cite[Lemma~2.3.1]{ngo1999faisceaux} and Mirkovi\'c--Vybornov \cite[Theorem~3.1 and \S3.3]{mirkovic2003quiver}.
Cautis and Kamnitzer \cite{cautis2018quantum} studied the equivariant $K$-theory of $\SX_{\Delta^+}$, describing its $K$-ring via the convolution tower and constructing an affine Hecke algebra action from a smooth divisor and a lattice-forgetting fibration.
They compared the resulting correspondences with the Steinberg correspondences on the open cotangent bundle $T^*\CB$ \cite[\S11]{cautis2018quantum}.

In this paper, we study the integral equivariant cohomology rings of $\SX_\Psi$ for arbitrary root ideals $\Psi\subset\Delta^+$ and the operators defined by geometric correspondences, retaining the loop rotation action of $S=\C^\times$ throughout. We give a uniform operator presentation of $H^*_K(\SX_\Psi;\Z)$ for $K=T\times S$ and $G\times S$, and construct geometric realizations of the Demazure--Lusztig type operators for all simple reflections preserving $\Psi\cup(-\Psi)$.

To state our results precisely, fix $n\geq2$ and set $[n]=\{1,\dots,n\}$. Let $S=\C^\times$ act on $\C(\!(t)\!)^n$ by $v(t)\mapsto v(at)$ for $a\in S$. Root ideals $\Psi\subset\Delta^+$ are in bijection with \emph{reverse Hessenberg functions}, namely weakly increasing maps $\rfunc:[n]\to\{0,\dots,n-1\}$ with $\rfunc(k)<k$ for all $k\in[n]$, via
\[
 \Psi_{\rfunc}=\{\epsilon_a-\epsilon_b\in\Delta^+\mid a\leq\rfunc(b)\},
\]
where $\epsilon_a$ is the $a$-th coordinate character of $T$. In the lattice model, $\SX_{\rfunc}=\SX_{\Psi_{\rfunc}}$ parametrizes flags $L_0\subset L_1\subset\cdots\subset L_n$ of $\C\llbracket t\rrbracket$-lattices in $\C(\!(t)\!)^n$ satisfying
\[
 L_0=t\C\llbracket t\rrbracket^n,\qquad
 \dim_\C(L_k/L_{k-1})=1,\qquad
 tL_k\subset L_{\rfunc(k)}\quad(1\leq k\leq n).
\]
Set
\[
 \Lambda=\Lambda_{T\times S}=\Z[t_1,\dots,t_n,\epsilon],\qquad
 R=\Lambda[x_1,\dots,x_n],\qquad
 f(z)=\prod_{a=1}^n(z+t_a),
\]
where $t_1,\dots,t_n\in H^2_T(\pt;\Z)$ are the equivariant weights of the standard representation $\C^n$, $\epsilon=c_1^S(\C_1)$ for the one-dimensional representation $\C_1$ of $S$ of weight $1$, and $\deg t_a=\deg\epsilon=\deg x_k=2$. Let $\tau_i$ be the $\Lambda$-algebra involution of $R$ swapping $x_i$ and $x_{i+1}$, and define the divided difference $\partial_i$ and the Demazure--Lusztig type operator $\BS_i$ on $R$ by
\[
 \partial_i g=\frac{g-\tau_i g}{x_i-x_{i+1}},\qquad
 \BS_i=\tau_i-\epsilon\partial_i
       =\id+(x_{i+1}-x_i-\epsilon)\partial_i
\]
(see \cite[\S12]{ginzburg1998geometric}).

For $1\leq j\leq n$, let $c_j=c_j^G(\C^n)\in H_G^{2j}(\pt;\Z)$ denote the equivariant Chern classes of the standard representation of $G$, also viewed as classes for $G\times S$ by pullback. Under restriction to $T\times S$, these classes map to the elementary symmetric polynomials $e_j(t_1,\dots,t_n)$.

\begin{theorem}[Theorem~\ref{thm:equiv_coho_xdelta} and Corollary~\ref{cor:equiv_coho_xdelta_G}]\label{thm:intro_ring}
For any reverse Hessenberg function $\rfunc$, there is a canonical isomorphism of graded $\Lambda$-algebras
\[
 H^*_{T\times S}(\SX_{\rfunc};\Z)\cong R/I_{\rfunc},\qquad
 I_{\rfunc}=(D_1(f),\dots,D_n(f)),
\]
where, with empty operator products understood as the identity,
\[
 D_k(f)=\partial_{k-1}\cdots\partial_{\rfunc(k)+1}
             \BS_{\rfunc(k)}\cdots\BS_1 f(x_1).
\]
Replacing $\Lambda$ by $\Z[c_1,\dots,c_n,\epsilon]$ and $f(z)$ by $z^n+c_1z^{n-1}+\cdots+c_n$ yields the presentation of $H^*_{G\times S}(\SX_{\rfunc};\Z)$. Under these isomorphisms for $K=T\times S$ and $G\times S$, $x_k$ maps to $-c_1^K(L_k/L_{k-1})$.
\end{theorem}

The proof of Theorem~\ref{thm:intro_ring} uses a polynomial recurrence to identify the reversed polynomial associated with $D_k(f)$, modulo $(D_1(f),\dots,D_{k-1}(f))$, with the Chern polynomial of $t^{-1}L_{\rfunc(k)}/L_{k-1}$. Since $D_k(f)\in\Lambda[x_1,\dots,x_k]$ is monic of degree $d_k=n-k+1+\rfunc(k)$ in $x_k$, the relations are triangular; successive division yields normal forms in the monomial basis $\{x_1^{e_1}\cdots x_n^{e_n}\mid 0\leq e_k<d_k\}$. Equivariant integration then extracts the coefficient of $x_1^{d_1-1}\cdots x_n^{d_n-1}$ (Corollary~\ref{cor:normal_forms}).

Next, let $S'$ be the set of simple reflections $s_i=(i,i+1)$ in the Weyl group $\mathfrak S_n$ preserving $\Psi\cup(-\Psi)$, and let $\Hdeg'$ be the subalgebra of the degenerate affine Hecke algebra $\Hdeg$ generated by $S'$ and $\Z[x_1,\dots,x_n,\epsilon]$. For $K=T\times S$ or $G\times S$, write $\Phi_K:R_K=\Lambda_K[x_1,\dots,x_n]\twoheadrightarrow H_K^*(\SX_\Psi;\Z)$ for the presentation map.

\begin{theorem}[Theorem~\ref{thm:geometric_hecke} and Corollary~\ref{cor:deg_hecke_action_xdelta}]\label{thm:intro_hecke}
For each $s_i\in S'$, define the degree-zero endomorphism $\mathcal S_i$ of $H_K^*(\SX_\Psi;\Z)$ as follows:
\begin{enumerate}
\item If $\alpha_i=\epsilon_i-\epsilon_{i+1}\in\Psi$, let $j_i:Y_i\hookrightarrow\SX_\Psi$ be the smooth Cartier divisor defined by the condition $tL_{i+1}\subset L_{i-1}$, and let $\pi_i:Y_i\to B_i$ be the morphism to the corresponding partial lattice flag variety $B_i$ obtained by forgetting $L_i$. Then $\pi_i$ is an equivariant $\PP^1$-bundle, and we set
\[
 \mathcal S_i=\id+j_{i*}\pi_i^*\pi_{i*}j_i^*.
\]
\item If $\alpha_i\notin\Psi$, then forgetting $L_i$ defines a global equivariant $\PP^1$-bundle $\pi_i:\SX_\Psi\to B_i$, and we set
\[
 \mathcal S_i=\id+c_1^K(\Omega^1_{\SX_\Psi/B_i}\otimes\C_{-1})\cup\pi_i^*\pi_{i*},
\]
where $\C_{-1}$ is the one-dimensional representation on which $G$ acts trivially and $S$ acts with weight $-1$.
\end{enumerate}
In both cases, $\mathcal S_i\Phi_K(g)=\Phi_K(\BS_i g)$ for all $g\in R_K$. Consequently, $\ker\Phi_K$ is stable under $\BS_i$, and the operators $\mathcal S_i$ ($s_i\in S'$), together with multiplication by $x_j$ and $\epsilon$, define a graded representation of $\Hdeg'$ on $H_K^*(\SX_\Psi;\Z)$.
\end{theorem}

Restriction to the open subbundle $T^*_\Psi\CB$ defines a surjective $\Hdeg'$-module homomorphism $H_K^*(\SX_\Psi;\Z)\twoheadrightarrow H_K^*(T^*_\Psi\CB;\Z)$ sending each $x_i$ to the corresponding tautological class on $T^*_\Psi\CB$ (also denoted by $x_i$). For $G\times S$, the target is the polynomial representation $\Z[x_1,\dots,x_n,\epsilon]$. This proves faithfulness over $\Z[\epsilon]$ for $G\times S$; restriction to $T\times S$ gives faithfulness for that group as well (Corollary~\ref{cor:faithful_hecke}). Equivariant duality then gives an embedding of $\Hdeg'$ into the convolution algebra of correspondence classes on $\SX_\Psi\times\SX_\Psi$ (Corollary~\ref{thm:faithful_convolution_embedding}). For $\Psi=\Delta^+$, the fiber product $\overline Z_i=Y_i\times_{B_i}Y_i$ is the closure of the simple Steinberg component $Z_i\subset T^*\CB\times_{\CN}T^*\CB$, relating this action to the classical Steinberg correspondences (Subsection~\ref{sec:steinberg}).

\paragraph{Organization of the paper.}
Section~\ref{sec:deg_hecke} establishes the polynomial recurrences for $D_k(f)$, and Section~\ref{sec_xpsi} proves Theorem~\ref{thm:intro_ring}.
Section~\ref{sec:pushpull} computes the Gysin pushforwards for inclusions $\SX_{\rfunc}\hookrightarrow\SX_{\rfunc'}$ (Proposition~\ref{prop:general_gysin}) and the push--pull operators for lattice-forgetting morphisms, and Section~\ref{sec:geometric_hecke} proves Theorem~\ref{thm:intro_hecke}.
Section~\ref{sec_TX_incl} treats open restriction, faithfulness, and the comparison with Steinberg correspondences. Section~\ref{sec:examples} works out explicit examples in ranks two and three.
